% Q1-Q2 问题重述；经 ACCT 于 2026-09-11 审核确认。
\section{问题重述}

\subsection{问题背景}

微网由用电负荷、光伏发电单元、储能装置及外部电网组成。光伏出力与负荷需求在日内并不完全匹配，储能可通过充放电实现电能在不同时段之间的转移，外部电网则为微网提供补充电能。因此，需要结合日内电价与负荷、光伏信息，统筹安排外网购电和储能运行，在满足供电需求的同时降低购电费用。对于连续多日运行，还需考虑预测偏差、紧急购电以及储能状态的跨日传递，使日前计划与日内实际执行相互衔接。

\subsection{问题一的要求}

\textbf{问题一：}根据给定的日内电价、负荷预测功率和光伏预测功率，将一天按附件原始顺序划分为 144 个 10 分钟时段，制定日前外网购电与储能充放电计划。在满足各时段能量平衡、储能充放电功率限制、储电量上下界以及首末储电量一致等条件下，使全天计划购电费用最小。给出逐时段购电量，汇总指定时段购电量、全天购电量与购电费用，并列出六个四小时区间的充放电量及日初、日末储电量，按附件模板填报。

\subsection{问题二的要求}

\textbf{问题二：}在问题一单日调度的基础上，进一步考虑全年运行中负荷与光伏预测存在偏差的情形。每天 0:00 应依据当时可获得的信息制定当日普通购电计划；进入日内运行后，根据实际负荷、实际光伏出力及当前储能状态调整实际取电与储能充放电。当计划购电、光伏和储能仍不能满足负荷需求时，以当时电价的 5 倍进行紧急购电。要求在保持储能状态跨日连续、满足设备运行约束且不利用未来实际信息的条件下，形成日前计划与日内执行相联动的购电策略，降低实际购电总成本，并按附件模板给出计划购电量、充放电量、日初与日末储电量以及紧急购电记录。
